\chapter{<Weiteres Kapitel>}
\label{chap:weiteres-kapitel}


Natürlich sollten Sie den Überschriften sinnvoll wählen und nicht "`<Weiteres Kapitel>"' schreiben. Auch sollten Sie alle "`<"' und "`>"' in den Überschriften entfernen.

Dieses Kapitel ist das Haupt-Kapitel Ihrer Arbeit. Typischerweise gibt es davon auch mehrere. Haben Sie z.B. einen Algorithmus, Verfahren o.ä. entwickelt, implementiert und getestet, macht es vielleicht Sinn, in  Kapitel 3 den Algorithmus konzeptuell und unabhängig von der konkreten Programmierung zu beschreiben; welche Lösungsalternativen haben Sie ggf. abgewogen? In Kapitel 4 könnten Sie dann implementierungsspezifische Details des Algorithmus beschrieben, z.B. welche Programmiermuster Sie eingesetzt haben. In Kapitel 5 könnten Sie dann die Evaluation des Algorithmus und seiner Implementierung beschreiben.

Ist Ihre Arbeit über eine empirische Studie, könnten Sie in Kapitel 3 die Planung der Studie beschreiben, warum Sie denken, dass diese Studie eine bestimmte Hypothese oder Fragestellung möglichst gut beantworten wird. In Kapitel 4 könnten Sie dann die Ergebnisse der Studie präsentieren und interpretieren. In Kapitel 5 können Sie über die methodische Herangehensweise reflektieren.

Details dazu besprechen Sie am besten mit Ihrer Betreuerin oder Ihrem Betreuer.